%% kapa: exam=1 round=24 year=2013 period=2 form=A subject=appraisal_law subject_ko=부동산관계법규 seq=21 number=61
%% source: Q-Net 감정평가사 기출문제 (cst00309) artlSeq=5073802 file=extracted/1008/5073802/2013년 제24회 감정평가사 1차시험 기출문제/2교시(A형).hwp
%% conversion: kordoc 4.24.51 -> tools/gen1.py (auto); review=proofread-sonnet; flags=none
\begin{question}{61}
부동산 가격공시 및 감정평가에 관한 법령상 요구되는 최소 인원수에 관한 설명으로 옳지 않은 것은? (단, 감정평가사는 자격등록 및 갱신등록의 거부사유가 없는 자임을 전제로 한다.)
\begin{choices}
  \item 감정평가법인에는 10인 이상의 감정평가사를 두어야 한다.
  \item 감정평가법인의 주사무소에 주재하는 최소 감정평가사의 수는 3명이다.
  \item 감정평가법인의 분사무소에 주재하는 최소 감정평가사의 수는 2명이다.
  \item 감정평가사합동사무소에 두는 감정평가사의 수는 5인 이상으로 한다.
  \item 국토교통부장관으로부터 업무의 위탁을 받을 수 있는 감정평가법인은 소속 감정평가사의 수가 50인 이상이어야 한다.
\end{choices}
\end{question}
\answer{4}
